\SourceLeaf{381}%
to the parallel on which the radius of curvature of the meridian equals the radius of a circle whose circumference equals the entire elliptical meridian. With the adopted flattening of 0,00324, the latitude of this parallel is \(45^0\,3'\,29''2\). The geodetic measurements give 475121,06 metres for the entire arc of this parallel between the meridians of Milan (observatory) and Usson, while observations of the gunpowder signals gave \(6^0\,1'\,41''7\) for the difference in longitude. These figures may be regarded as the principal result of the survey. Comparing such an arc of longitude with the result of a survey of an arc of latitude can in theory determine the Earth's flattening. By comparing their result in this way with the arc from Greenwich to Formentera, the authors find a flattening of \(\frac{1}{241}\). How little reliance can be placed on this method is most strikingly apparent when the individual sections are treated in the same way as the whole arc. Comparing the same meridian arc with the section d’Usson --- Colombier, the reviewer thus finds a flattening of \(\frac{1}{169}\), with the second section Colombier --- Mont Cenis \(\frac{1}{226}\), with the fourth, Turin --- Milan, \(\frac{1}{110}\), and with the third, Mont Cenis --- Turin, instead an elongation of \(\frac{1}{82}\). So there is little to be expected from this in determining the Earth's dimensions more precisely. All the more important are the results as fresh confirmation of the irregularity of the Earth's figure, which is most pronounced precisely in northern Italy. These irregularities stand out most clearly when the astronomically determined differences in longitude are compared with those calculated from the geodetic measurements on a plausible hypothesis about the overall figure of the Earth. The authors made this calculation using a flattening of 0,00324 and an equatorial radius of 6376986 metres. The resulting differences in longitude west of Milan, expressed in time, are
\begin{center}
\begin{tabular}{l|r@{\quad}l|r@{\quad}l|r}
&\multicolumn{2}{c|}{astronomical}&\multicolumn{2}{c|}{geodetic}&Difference\\
Turin & \(5'\) & \(58''85\) & \(6'\) & \(0''93\) & \(-2''08\)\\
Mont Cenis & \(9\) & \(0,20\) & \(8\) & \(59,49\) & \(+0,71\)\\
Colombier & \(13\) & \(44,23\) & \(13\) & \(43,84\) & \(+0,39\)\\
D’Usson & \(24\) & \(6,78\) & \(24\) & \(8,02\) & \(-1,24\)
\end{tabular}
\end{center}
